% I. Литературен преглед — ЗАДЪЛЖИТЕЛЕН за ОКС „Магистър“ (REQUIREMENTS.md А.3).
% Ориентировъчен обем: 8–10 стр.
\section{Анализ на предметната област и литературен преглед}
\label{sec:literature}

Разделът е подреден от общото към частното: първо класификация на подходите според
това какво в крайна сметка рисува пиксел на екрана, после конкретните представители на
всяко семейство, и накрая --- по-подробно --- архитектурата на реализацията,
която в тази работа изпълнява ролята на еталонен измервателен инструмент, а не на образец
за копиране. Последното е разгледано подробно, защото границата, която тя въвежда, е и
предметът на измерването в Раздел~\ref{sec:results}.

\subsection{Подходи за разработка на многоплатформени графични интерфейси}
\label{subsec:approaches}

Съществуващите решения се разделят на три семейства според това \emph{какво} в крайна
сметка рисува пиксел на екрана. Разделението не е самоцелна класификация --- то
определя всичко останало: вида на крайния контрол, размера на артефакта, достъпа до
платформените услуги и --- което е предметът на настоящата работа --- \emph{дали изобщо
съществува понятие „правилен вид“}, спрямо което решението да бъде проверено.

\subsubsection{Абстракция над native контроли}

Едно описание се преобразува към контролите на \emph{самата} платформа: списъкът е
native списък, бутонът е native бутон. Предимствата идват от само себе си ---
външният вид, поведението при въвеждане, достъпността и езиковите настройки идват от
платформата и се променят заедно с нея.

Цената е, че абстракцията трябва да \emph{съгласува} инструментариуми, които се различават
не само по имена, а по модел: различни правила за разположение, различни жизнени цикли,
различни конвенции за координати. Именно тук възниква въпросът, който настоящата работа
измерва --- доколко такова съгласуване е постижимо и къде се проваля.

\subsubsection{Собствен растеризатор}

Фреймуоркът не използва native контроли, а рисува всеки елемент сам, върху графично
платно. Резултатът е еднакъв на всички платформи \emph{по замисъл}, което премахва
целия клас проблеми на съгласуването.

Замяната обаче е реална: интерфейсът престава да бъде native. Той не наследява
автоматично промените във външния вид на платформата, а достъпността, въвеждането с
клавиатура и текстовите услуги трябва да бъдат възпроизведени във фреймуорка. Артефактът
носи и самия растеризатор.

\subsubsection{Вграден уеб изглед}

Интерфейсът е уеб документ в native обвивка. Този подход има най-ниска цена на входа и
най-голяма отдалеченост от платформата; той е и подходът с най-обилно изследвана цена в
литературата~\cite{biornhansen2020overhead}.

\subsubsection{Критерии за съпоставяне}

\begin{table}[H]
\caption{Съпоставяне на трите подхода за многоплатформен графичен интерфейс}
\label{tab:approaches}
\centering\small
\begin{tabular}{@{}L{3.3cm}L{3.3cm}L{3.3cm}L{3.3cm}@{}}
\hline
\textbf{Критерий} & \textbf{Над native контроли} & \textbf{Собствен растеризатор} & \textbf{Вграден уеб изглед} \\
\hline
краен контрол & native & собствен & документен елемент \\
външен вид след промяна в платформата & следва я сам & изисква работа & изисква работа \\
достъпност, въвеждане & от платформата & във фреймуорка & от уеб слоя \\
размер на артефакта & малък & + растеризатор & + уеб механизъм \\
среда за изпълнение & по избор & обичайно да & да \\
\textbf{наличие на еталон} & \textbf{да} & не (сам си е еталон) & частично \\
\hline
\end{tabular}
\end{table}

Последният ред е решаващият за настоящата работа. При собствен растеризатор въпросът
„съответства ли изобразеното“ не стои --- изобразеното \emph{по определение} е
правилното, защото няма втора инстанция, спрямо която да се сравнява. Само при
абстракция над native контроли твърдението за съответствие е смислено, проверимо ---
и, както показва §\ref{subsec:lit_summary}, непроверено.

% \begin{table}[H]
% \caption{Сравнение на подходите за многоплатформени графични интерфейси}
% \label{tab:approaches}
% \end{table}

\subsection{Съществуващи многоплатформени фреймуорци}
\label{subsec:frameworks}

\subsubsection{Xamarin.Forms и .NET MAUI}

Историческата линия е поучителна, защото съдържа точно един архитектурен завой.
Xamarin.Forms~\cite{xamarin_forms_docs} свързва всяка контрола с платформената чрез \term{renderer} --- клас, който
наследява по йерархия и включва както създаването на native изгледа, така и
пренасянето на всяко свойство. Наследяването означава, че разширяването на една контрола
изисква наследяване на нейния renderer, а той е обвързан с платформата.

.NET MAUI~\cite{maui_docs,maui_handlers} заменя този механизъм с \term{handler}: обектът вече не се наследява, а носи
\emph{таблица} от съответствия „свойство $\rightarrow$ действие“. Разликата е между
разширяване чрез наследяване и разширяване чрез вписване в таблица --- второто е
съставимо, не изисква познаване на йерархията и позволява една контрола да бъде
доразвита, без нищо да се наследява. Именно този шаблон е предметът на
§\ref{subsec:maui_arch} и основата, върху която стъпва настоящата работа.

\subsubsection{Flutter}

Представител на семейството със собствен растеризатор~\cite{flutter_docs}. Изгражда интерфейса от собствени
примитиви върху графично платно, което дава пълна еднаквост между платформите и пълна
свобода в анимациите, но и означава, че всяка контрола, всяко поведение при въвеждане и
всяка услуга за достъпност съществуват във фреймуорка, а не в платформата.

По критериите от \ref{subsec:approaches} Flutter заема крайното положение на второто
семейство: пълна еднаквост между платформите, пълна независимост от промени в native
инструментариума, и артефакт, който носи целия слой за изобразяване. За настоящото
изследване той е показателен по обратен начин --- при него описаният тук метод за
измерване е \emph{неприложим}, защото липсва втора независима реализация на същото
описание, спрямо която да се сравнява.

\subsubsection{Qt}

Най-старият зрял C++ фреймуорк в областта~\cite{qt_docs} и единственият пряк съсед на настоящата работа
по език. Съдържа и двата подхода: по-старият слой със \term{widgets} рисува сам, като
подражава на външния вид на платформата, а по-новият декларативен слой е самостоятелен
растеризатор със собствен език за описание.

Разглеждането му е особено уместно, защото Qt решава --- на същия език --- и задачата за
собствеността без събирач на боклук: дървото от обекти е с притежаване от родител към
дете, а наблюдаващите препратки са отделен тип. Настоящата работа стига до сходно
разпределение на ролите по независим път (§\ref{subsec:ownership}), което е довод, че
решението е принудено от задачата, а не въпрос на предпочитание.

\subsubsection{Avalonia, Slint и React Native}

\textbf{Avalonia}~\cite{avalonia_docs} е собствен растеризатор върху управлявана среда, с декларативен
език, силно повлиян от предходно поколение настолни технологии; позиционира се като
преносим наследник на тези технологии, което го поставя във второто семейство въпреки
близостта на езика му за описание до първото.

\textbf{Slint}~\cite{slint_docs} също рисува сам, но с вграден собствен език за описание и с изричен
прицел към устройства с ограничени ресурси --- което го прави най-близкия по цели
представител на второто семейство за вградени приложения.

\textbf{React Native}~\cite{reactnative_docs} принадлежи на първото семейство: изгражда native контроли, но
описанието се изпълнява в скриптов език, а границата между двете страни е асинхронна.
Точно тази граница е предметът на измерването
в~\cite{biornhansen2020overhead} --- изследване, което мери \emph{цената} на връзката, но
не и съответствието на изобразеното.

\subsection{Архитектурата на .NET MAUI}
\label{subsec:maui_arch}

Централната идея, върху която стъпва цялата работа, е че разглежданата реализация
\textbf{не е растеризатор}. Тя не рисува нито един пиксел сама. Тя е абстракция, която
преобразува едно многоплатформено описание на интерфейса към няколко native
инструментариума, и е организирана около един-единствен повтарящ се шаблон, приложен
еднакво за всяка контрола.

\subsubsection{Шаблонът Handler: виртуален изглед, mapper, native контрола}
Връзката е тристранна и е показана на \figref{fig:handler}. Най-простият ѝ срез ---
бутон --- се състои от: \emph{договор} (интерфейс само със свойства, без изобразяване);
\emph{обработчик}, който носи две таблици --- за свойства и за императивни операции;
и \emph{native контрола} на всяка платформа. Таблицата за свойства съпоставя име на
свойство с функция, която прилага текущата му стойност; таблицата за операции прави
същото за действия, които нямат стойност (например „получи фокус“).

Границата минава точно тук: договорът и таблиците са преносими, а само реализациите на отделните
действия са платформени. Затова добавянето на инструментариум е \emph{добавяне} на
реализации, а не изменение на съществуващи --- свойство, което Раздел~\ref{sec:design}
запазва съзнателно.

\begin{figure}[H]
\centering
\begin{tikzpicture}[font=\footnotesize,
  b/.style={draw, rounded corners=2pt, align=center, inner sep=4pt, minimum height=0.75cm},
  nat/.style={b, fill=red!8, text width=2.5cm},
  ar/.style={-Latex, thick}]

  \node[b, fill=blue!6, text width=3.1cm] (vv) {\textbf{Виртуален изглед}\\\emph{договор}, без рисуване};
  \node[b, fill=green!10, text width=3.6cm, right=1.5cm of vv] (h)
      {\textbf{Обслужващ обект}\\таблица „свойство $\rightarrow$ действие“\\+ таблица за операции};

  \node[nat, right=1.9cm of h, yshift=1.55cm] (n1) {UIKit};
  \node[nat, right=1.9cm of h, yshift=0.75cm] (n2) {AppKit};
  \node[nat, right=1.9cm of h, yshift=-0.05cm] (n3) {Android View};
  \node[nat, right=1.9cm of h, yshift=-0.85cm] (n4) {WinUI};
  \node[nat, right=1.9cm of h, yshift=-1.65cm, fill=gray!10] (n5) {без екран};

  \draw[ar] (vv) -- node[above, font=\scriptsize] {промяна} (h);
  \draw[ar] (h.east) -- (n1.west);
  \draw[ar] (h.east) -- (n2.west);
  \draw[ar] (h.east) -- (n3.west);
  \draw[ar] (h.east) -- (n4.west);
  \draw[ar] (h.east) -- (n5.west);

  \node[below=0.5cm of h, align=center, font=\scriptsize, text width=8cm]
     {един договор $\rightarrow$ една таблица $\rightarrow$ $N$ инструментариума};
\end{tikzpicture}
\caption{Шаблонът на преобразуването. Виртуалният изглед носи само свойства; обработчикът съпоставя всяко свойство на действие, което го прилага върху native контролата.
Смяната на инструментариума заменя само дясната колона.}
\label{fig:handler}
\end{figure}

\subsubsection{Слоеста организация и ред на зависимостите}
Зависимостите между слоевете вървят в една посока --- графични примитиви, договори,
обработчици, платформени слоеве, разположение, контроли, декларативен слой, услуги
на устройството, инициализация --- и този ред е задължителен и при независимо изграждане
на такъв слой. Причината не е организационна: слой, който може да достигне до
платформения, престава да бъде преносим, дори ако в момента не го достига. Пълното
разглеждане на структурата, вече в реализирания вид, е в §\ref{subsec:structure}.

\subsubsection{Системата от свойства и свързването на данни}
Системата от свойства е вторият механизъм, който се повтаря навсякъде. Всяко свойство
е описано от споделен описател (име, стойност по подразбиране, обратни извиквания), а
всяка инстанция държи само стойността. Върху това стъпват три неща: известяване при
промяна --- което задвижва и таблиците от §\ref{subsec:maui_arch}; \term{привеждане}
(coercion) на стойността към допустим обхват; и \term{предимство} между едновременно
валидни източници на една и съща стойност --- подразбираща се, от стил, от свързване, от
ръчно присвояване и т.н.

Предимството е частта, която най-лесно се подценява. То не е вътрешна подробност, а
наблюдаемо поведение: то определя дали изрично зададен от приложението цвят надделява над
цвят от стил, и дали „незададено“ се различава от „зададено на стойността по
подразбиране“. Раздел~\ref{sec:results} показва, че точно това разграничение е източник
на цял клас разминавания.

\subsection{Изводи от литературния преглед}
\label{subsec:lit_summary}

Прегледът води до четири констатации, всяка от които се връзва със задача от Увода.

\textbf{Първо.} Твърдението за еднакво поведение върху различни платформи е \emph{общо
за цялото семейство} решения, които стъпват на native контроли, и никъде не е придружено
от количествена проверка. Съпоставителните изследвания измерват друго --- режийни разходи
на границата към платформата~\cite{biornhansen2020overhead}, потребление на ресурси,
удобство за разработчика --- но не и степента, в която двете реализации на едно описание
изобразяват едно и също.

\textbf{Второ.} Причината за тази празнина не е незаинтересованост, а липса на
\term{оракул}: за разлика от изчислителния код, „правилният“ вид на интерфейс не е
записан във вид, годен за машинно сравнение~\cite{barr2015oracle}. Изследванията върху
визуално тестване на интерфейс~\cite{garousi2017vgt} заобикалят това, като сравняват
система със \emph{самата нея} в предишен момент --- което открива регресии, но не може да
провери твърдение за съответствие между две различни реализации.

\textbf{Трето.} Съществува обаче положение, при което оракул е наличен: когато една зряла
реализация се приеме за меродавна и втора, независима реализация изпълнява \emph{същото}
описание. Тогава разликата между двете изображения е измерима величина. Това е
наблюдението, върху което стъпва настоящата работа.

\textbf{Четвърто.} Архитектурата с таблично преобразуване (§\ref{subsec:maui_arch}) е
особено подходяща за такава проверка, защото границата между преносимото и платформеното е
\emph{един} и е явен. Реализация, която го запазва, позволява разминаванията да бъдат
отнасяни към определен слой; реализация, която го събаря в един растеризатор, не
позволява това.

Тези четири констатации определят и подредбата на останалото изложение: избор на
средства и стратегия (Раздел~\ref{sec:alternatives}), запазване на тази граница при проектиране
и реализация (Раздели~\ref{sec:design} и~\ref{sec:implementation}), изграждане на оракула
в работеща методика (Раздел~\ref{sec:method}) и получените от нея резултати
(Раздел~\ref{sec:results}).
